\documentclass[10pt,a4paper]{amsart}
\usepackage[margin=19mm]{geometry}
\usepackage{amsmath,amssymb,mathtools}
\usepackage[scr=rsfs,cal=euler]{mathalfa}
\usepackage{xfrac}
\usepackage[unicode,hidelinks,bookmarks=false]{hyperref}
\newcommand{\ie}{i.e.\ }
\newcommand{\cf}{cf.\ }
\newcommand{\resp}{respectively\ }
\newcommand{\Subsection}{\subsection}
\providecommand{\nobreakdash}{\nobreak\hbox{-}}
\setlength{\emergencystretch}{2em}
\multlinegap0pt
\usepackage{bm}
\usepackage{bbm}
\usepackage{dsfont}
\usepackage{xspace}
\usepackage{cancel}
\usepackage{mathtools}
\usepackage{amsthm}
\makeatletter
\@ifundefined{thm}{\newtheorem{thm}{Thm}}{}
\@ifundefined{lem}{\newtheorem{lem}[thm]{Lemma}}{}
\makeatother
\def\Cay{{\sf Cay}} \def\Ker{{\sf Ker}}
\title[Finite groups with quadratic splitting fields for all Cayley]{Finite groups with quadratic splitting fields for all Cayley graphs}
\author{Majid Arezoomand, Alireza Abdollahi, Tao Feng, Zeinab Akhlaghi}
\date{Source: arXiv:2607.02973v1 (CC BY 4.0)}
\begin{document}
\maketitle

\noindent\emph{Source and licence.} Finite groups with quadratic splitting fields for all Cayley graphs. Version arXiv:2607.02973v1 (CC BY 4.0), distributed under the Creative Commons Attribution 4.0 International licence, as recorded in the arXiv licence metadata for this exact version and shown on the abstract page https://arxiv.org/abs/2607.02973v1. Full rights notice: licenses/arxiv-2607.02973-CC-BY-4.0.txt.\par\smallskip

\noindent\textit{Editorial scope.} Counted unit: the Lem whose source label is \texttt{32yes} in the article (source0.tex lines 478--487, source label \texttt{32yes}), delivered together with the article's own complete original proof of it (source0.tex lines 488--583, one \texttt{proof} environment of 96 lines). No mathematics is written, repaired or re-derived: statement and proof are the article's own text line for line. Required context retained uncounted: none. Every definition, display and label of those lines is retained unchanged. Omitted from this package and not redistributed: 1--477, 584--1922. Nothing inside a retained excerpt is omitted or altered: every retained body is delivered exactly as the article's lines are written, so no omission receipt of any kind accompanies this package. Citations: the retained excerpts carry no citation command at all, so no bibliography entry of the article is transcribed and no reference is redistributed. Reference adaptation: none was needed and none was made; every reference inside the delivered unit resolves inside it, so no unresolved reference or citation is produced. Printed slips kept literal: no printed word, symbol, hypothesis or number of the counted statement or of its proof was altered. Layout and class: the article's own class is \texttt{amsart} with packages including amsmath, amsthm, amsfonts, amssymb, latexsym, mathrsfs, graphicx, hyperref, enumerate, enumitem, color, tikz, ytableau; the wrapper is the lane's pinned \texttt{amsart} editorial wrapper at 10pt on A4, which restates the theorem-like environments the retained text needs and adds the article's own macro definitions that the retained text needs (\texttt{\textbackslash Cay}), with extra wrapper packages (bm, bbm, dsfont, xspace, cancel, mathtools). The delivered numbering is the unit's own. Assets: no retained span contains a figure environment, a graphics-inclusion command, a PDF-page inclusion, a table or a TikZ picture; no image file is copied into this package, the package declares no asset dependency and no image file is redistributed with it. Licence: version arXiv:2607.02973v1 of this article is distributed under the Creative Commons Attribution 4.0 International licence, as recorded in the arXiv licence metadata for this exact version and shown on the abstract page https://arxiv.org/abs/2607.02973v1. A full rights notice is delivered at licenses/arxiv-2607.02973-CC-BY-4.0.txt. Characters: every retained excerpt is pure ASCII, so the delivered engine, XeLaTeX with its default Latin Modern text font, reports no missing character.
\begin{lem}\label{32yes}
	The following non-abelian groups of order 32 are Cayley $2$-integral:
	\begin{itemize}
		\item[(1)] $G_{32,50}=\langle a,b,c,d\mid a^2=b^4=d^2=c^2b^2=cc^{b^{-1}}=(ab)^2=dd^b=a^ca=dd^c=a^db^2a=1\rangle$,
		\item[(2)] $G_{32,35}=\langle a,b,c\mid b^4=c^4=1, a^2=b^{-2}, ab=b^{-1}a,ac=c^{-1}a,bc=cb\rangle$,
		\item[(3)] $G_{32,26}=\langle a,b,c\mid  b^4=a^4=c^4=1, ab=b^{-1}a, ca=ac, cb=bc\rangle\cong\Bbb Z_4\times Q_8$,
		\item[(4)] $G_{32,23}=\langle a,b,c\mid c^2=b^4=a^4=b^ab=cc^a= cc^b=a^2b^{-1}a^2b=(b^{-1}a^{-2}b^{-1})^2=1 \rangle$.
		%	\item[(5)] ????
	\end{itemize}
\end{lem}

\begin{proof}
	$(1)$ Let $S$ be an inverse-closed subset of $G=G_{32,50}$ and $1\notin S$. Then
	$S$ is a union of the following sets
	
	
	$$\{ a \}, \{ b^2 \}, \{ d \}, \{ ab \}, \{ ab^2 \},\{ db^2 \},	\{ abc \},\{ ab^3 \},\{ acd \}, \{ abcb^2 \}, \{ acdb^2 \},$$
	$$ 
	\{ b, b^3\}, \{ c, cb^2 \},\{ ac, acb^2 \},
	\{ ad, adb^2 \}, \{ bc, bcb^2 \}, \{ bd, bdb^2 \},$$
	$$\{ cd, cdb^2 \}, \{ abd, abdb^2 \},   \{ bcd, bcdb^2 \},
	\{ abcd, abcdb^2 \}. $$
	By GAP, we see that $G$ has exactly $17$ irreducible representations, one of them, say $\varphi$, has degree $4$ and the others have degrees $1$ and are integer-valued. We claim that the eigenvalues of $\varphi(S)$ are in $\Bbb Q[\sqrt{m}]$ for some fixed positive integer $m$, which  proves $\Cay(G,S)$ is Cayley $2$-integral.  By GAP, we have
	\begin{equation*}
		\varphi(a)=\begin{bmatrix}
			0& 0&1&0 \\
			0&0&0&1\\
			1&0&0&0\\
			0&1&0&0
		\end{bmatrix},~~	\varphi(b)=\begin{bmatrix}
			0& 1&0&0 \\
			-1&0&0&0\\
			0&0&0&-1\\
			0&0&1&0
		\end{bmatrix},
	\end{equation*}
	\begin{equation*}
		\varphi(c)=\begin{bmatrix}
			-\textbf{i}& 0&0&0 \\
			0&\textbf{i}&0&0\\
			0&0&-\textbf{i}&0\\
			0&0&0&\textbf{i}
		\end{bmatrix},~~	\varphi(d)=\begin{bmatrix}
			-1& 0&0&0 \\
			0&-1&0&0\\
			0&0&1&0\\
			0&0&0&1
		\end{bmatrix},
	\end{equation*}
	where $\textbf{i}=\sqrt{-1}$. Then $\varphi(b^2)=-I$, where $I$ is the $4\times 4$ identity matrix. Hence if $X$ is a set of size $2$ chosen from the above sets, then $\varphi(X)=O$, where $O$ is the $4\times 4$ zero matrix. Let $S=S_1\cup S_2$, where $S_j$ is a union of the above sets of size $j$, $j=1,2$. Then $\varphi(S)=\varphi(S_1)+\varphi(S_2)=\varphi(S_1)+O=\varphi(S_1)$. Since $\varphi(b^2)=-I$, we may assume that $b^2\notin S_1$. Hence $S_1$ is a union of the sets 
	$$\{ a \},  \{ ab^2 \},\{ d \}, \{ db^2 \}, \{ ab \}, \{ ab^3 \},	\{ abc \}, \{ abcb^2 \}, \{ acd \},  \{ acdb^2 \}.$$
	Since $\varphi(x)+\varphi(xb^2)=O$ for any $x\in G$, we may assume that $S_1$ has at most $5$ elements and if $x\in S_1$ then $xb^2\notin S_1$. By considering all possibilities for $S$ and by a tedious computation, one can see that the characteristic polynomial of $\varphi(S)$ is one of the following
	\[ (x-1)^2(x+1)^2,~ (x^2-2)^2,~(x^2-3)^2, (x-2)^2(x+2)^2,~(x^2-5)^2,\]
	which proves our claim and completes the proof.
	
	$(2)$. One can check that $Z(G)=\{1,a^2,c^2,a^2c^2\}$.  Furthermore, the only non-identity elements of $G$ of order $2$ are $a^2,c^2,a^2c^2$.  Let $S$ be an inverse-closed subset of $G=G_{32,35}$ and $1\notin S$. Then
	$S$ is a union of the following sets
	$$\{a^2\},~\{c^2\},~\{a^2c^2\},$$
	$$
	\{a, a^3\},
	\{b, ba^2\},
	\{ c, c^3\},
	\{ab, aba^2\},
	\{ac, aca^2\},
	\{ ac^2, a^3c^2\},
	\{ bc, bca^2c^2\},$$
	$$
	\{bc^2, ba^2c^2\},
	\{ca^2, ca^2c^2\},
	\{abc, abca^2\},
	\{abc^2, aba^2c^2\},
	\{ac^3, ac^3a^2\},
	\{bca^2, bc^3\},
	\{abc^3, abc^3a^2\}.
	$$
	
	Clearly, if $\{x_1,x_2\}$ is a set of the above list, then $x_2=x_1y$, where $y\in\{a^2,c^2,a^2c^2\}$. Let $\rho$ be an irreducible representation of $G$ of degree $d$, and $S=X\cup Y$, where $X$ and $Y$ are a union of the sets of size $2$ and size $1$, chosen from the above list,  respectively. Then $\rho(S)=\rho(X)+\rho(Y)$. Since the elements of $Y$ are involutions, $\rho(Y)$ is an integer multiple of  $I$, where $I$ is the $d\times d$ identity matrix. So it is enough to find the eigenvalues of $\rho(X)$. By considering all possibilities for $X$ and all irreducible representations of $G$, one can see that, the characteristic polynomial of $\Cay(G,X)$ is of the form $f_1(x)f_2(x)$ where $f_1(x)$ is a product of factors of degree $1$, and either $f_2(x)=1$ or $f_2(x)$ is one of the followings:
	\[(x^2-8)^4,~(x^2-20)^2,~(x^2-32)^2,(x^2+4x-16)^2,~(x^2-4x-16)^2,\]
	\[(x^2-4x-28)^2,(x^2+4x-28)^2,~(x^2-4x-4)(x^2+4x-4).\]
	This implies that $G$ is Cayley $2$-integral.
	
	$(3)$ One can see that $Z(G)=\{1,c,b^2,a^2,cb^2,ca^2,b^2a^2,cb^2a^2\}$.
	Let $S$ be an inverse-closed subset of $G=G_{32,26}$ and $1\notin S$. Then
	$S$ is a union of the following sets
	\[\{a^2\},\{b^2\},\{a^2b^2\},\]
	\[\  \{a, a^3\}, \{b, b^3\} ,\{ c, cb^2a^2 \},\{ ab, aba^2 \},\{ ac, acb^2 \},
	\{ ab^2, ab^2a^2 \},\{ bc, bca^2 \},\]
	\[\{ ba^2, b^3a^2 \},\{ cb^2, ca^2 \},\{ abc, abcb^2 \},
	\{ ab^3, ab^3a^2 \},\{ a^3c, a^3cb^2 \},\{ b^3c, b^3ca^2 \},\{ a^3bc, a^3bcb^2\}.
	\]
	Clearly, if $\{x_1,x_2\}$ is a set of the above list, then $x_2=x_1y$, where $y\in\{a^2,b^2,a^2b^2\}$. Let $\rho$ be an irreducible representation of $G$ of degree $d$, and $S=X\cup Y$, where $X$ and $Y$ are a union of the sets of size $2$ and size $1$, chosen from the above list,  respectively. Then $\rho(S)=\rho(X)+\rho(Y)$. Since the elements of $Y$ are involutions, $\rho(Y)$ is an integer multiple of $I$, where $I$ is the $d\times d$ identity matrix. So it is enough to find the eigenvalues of $\rho(X)$. By considering all possibilities for $X$ and all irreducible representations of $G$, one can see that, the characteristic polynomial of $\Cay(G,X)$ is of the form $f_1(x)f_2(x)$ where $f_1(x)$ is a product of factors of degree $1$, and either $f_2(x)=1$ or $f_2(x)$ is one of the followings:
	\[(x^2-8)^4,~(x^2-12)^4.\]
	This implies that $G$ is Cayley $2$-integral.
	
	$(4)$
	One can see that $Z(G)=\{1,c,b^2,a^2,cb^2,ca^2,b^2a^2,cb^2a^2\}$.
	Let $S$ be an inverse-closed subset of $G=G_{32,23}$ and $1\notin S$. Then
	$S$ is a union of the following sets
	\[\{c\}, \{b^2\}, \{a^2\}, \{cb^2\}, \{ca^2\}, \{b^2a^2\}, \{cb^2a^2\},\]
	\[ \{ a, a^3 \}, \{ b, b^3 \}, \{ ab, aba^2 \}, \{ ac, aca^2 \}, \{ ab^2, ab^2a^2 \}, \{ bc, bcb^2 \},\]
	\[\{ ba^2, b^3a^2 \}, \{ abc, abca^2 \}, \{ ab^3, ab^3a^2 \},
	\{ acb^2, acb^2a^2 \}, \{ bca^2, bcb^2a^2 \}, \{ abcb^2, abcb^2a^2 \}.
	\]
	Clearly, if $\{x_1,x_2\}$ is a set of the above list, then $x_2=x_1y$, where $y\in\{a^2,b^2,a^2b^2\}$. Let $\rho$ be an irreducible representation of $G$ of degree $d$, and $S=X\cup Y$, where $X$ and $Y$ are a union of the sets of size $2$ and size $1$, chosen from the above list,  respectively. Then $\rho(S)=\rho(X)+\rho(Y)$. Since the elements of $Y$ are involutions, $\rho(Y)$ is an integer multiple of $I$, where $I$ is the $d\times d$ identity matrix. So it is enough to find the eigenvalues of $\rho(X)$. By considering all possibilities for $X$ and all irreducible representations of $G$, one can see that, the characteristic polynomial of $\Cay(G,X)$ is of the form $f_1(x)f_2(x)$ where $f_1(x)$ is a product of factors of degree $1$, and either $f_2(x)=1$ or $f_2(x)$ is one of the followings:
	\[(x^2-8)^4,~(x^2-20)^2,~(x^2-32)^2.\]
	This implies that $G$ is Cayley $2$-integral.
\end{proof}

\end{document}
